% Chapter 4, Topic _Linear Algebra_ Jim Hefferon
%  http://joshua.smcvt.edu/linearalgebra
%  2001-Jun-12
\topic{Rekurensi Linear}
\index{rekurensi linear|(}
\index{relasi rekurensi|(}

Pada tahun 1202, Leonardo dari Pisa, yang dikenal sebagai Fibonacci,
mengajukan masalah berikut.
\begin{quotation}
  \noindent Seorang pria menaruh sepasang kelinci di suatu tempat
  yang dikelilingi tembok pada semua sisinya.
  Berapa pasang kelinci dapat dihasilkan oleh pasangan itu dalam setahun
  jika diasumsikan bahwa setiap bulan tiap pasangan melahirkan satu pasangan
  baru yang mulai berkembang biak pada bulan keduanya?
\end{quotation}
Model ini mengembangkan model dasar pertumbuhan populasi secara eksponensial
dengan memperhitungkan bahwa kelinci yang baru lahir belum subur selama suatu
periode, dalam hal ini satu bulan.
Namun, model ini tetap mempertahankan asumsi penyederhanaan lain, misalnya
bahwa tidak ada usia yang setelahnya kelinci menjadi tidak subur.

Untuk memperoleh jumlah seluruh pasangan pada bulan berikutnya,
kita menjumlahkan banyak pasangan yang masih hidup ketika memasuki bulan itu
dengan banyak pasangan yang baru akan lahir pada bulan tersebut.
Besaran terakhir ini sama dengan banyak pasangan yang akan mampu
berkembang biak ketika memasuki bulan berikutnya, yaitu pasangan yang pada
bulan berikutnya sudah hidup setidaknya dua bulan.
\begin{equation*}
  F(n)=F(n-1)+F(n-2)  \qquad \text{dengan $F(0)=0$, $F(1)=1$}
  \tag{$*$}
\end{equation*}
Di ruas kiri terdapat suatu
\definend{relasi rekurensi}.\index{rekurensi}
Nama itu digunakan karena $F$ muncul kembali dalam persamaan yang
mendefinisikannya.
Di ruas kanan terdapat kondisi awal.
Dari~($*$), kita dapat menghitung $F(2)$, $F(3)$, dan seterusnya,
hingga memperoleh jawaban atas pertanyaan Fibonacci.
\begin{center}
  \begin{tabular}{r|ccccccccccccc}
    \textit{bulan}~$n$
     &0  &1  &2  &3  &4  &5  &6  &7  &8  &9  &10  &11  &12  \\ \hline
    \textit{pasangan}~$F(n)$
     &0  &1  &1  &2  &3  &5  &8  &13  &21  &34  &55  &89  &144   
  \end{tabular}
\end{center}
Kita akan menggunakan aljabar linear untuk memperoleh rumus yang menghitung
$F(n)$ tanpa harus lebih dahulu menghitung nilai-nilai antara $F(2)$, $F(3)$,
dan seterusnya.

Kita mulai dengan menyatakan ($*$) dalam bentuk matriks.
\begin{equation*}
  \colvec{F(n) \\ F(n-1)}
  =
  \begin{mat}[r]
    1  &1   \\
    1  &0
  \end{mat}
  \colvec{F(n-1) \\ F(n-2)}
  \qquad\text{dengan $\colvec{F(1) \\ F(0)}=\colvec{1 \\  0}$}
\end{equation*}  
Tuliskan $T$ untuk matriks tersebut dan $\vec{v}_{n}$ untuk vektor yang
komponennya $F(n)$ dan $F(n-1)$, sehingga
$\vec{v}_{n}=T^{n-1}\vec{v}_1$ untuk $n\geq 1$.
Jika kita mendiagonalkan $T$, kita memperoleh cara cepat untuk menghitung
pangkat-pangkatnya:~jika $T=PDP^{-1}$, maka $T^n=PD^nP^{-1}$, dan pangkat
ke-$n$ dari matriks diagonal $D$ adalah matriks diagonal yang entri-entrinya
merupakan pangkat ke-$n$ dari entri-entri $D$.

Persamaan karakteristik $T$ adalah $\lambda^2-\lambda-1=0$.
Rumus kuadrat memberikan akar-akarnya, yaitu $(1+\sqrt{5})/2$ dan
$(1-\sqrt{5})/2$.
(Keduanya kadang-kadang disebut ``rasio emas;''
lihat \cite{Falbo}.)
Pendiagonalan memberikan hasil berikut.
\begin{equation*}
  \begin{mat}[r]
    1  &1  \\
    1  &0
  \end{mat}
  =\begin{mat}
     \frac{1+\sqrt{5}}{2}  &\frac{1-\sqrt{5}}{2} \\
     % \sfrac{1+\sqrt{5}}{2}  &\sfrac{1-\sqrt{5}}{2} \\
     1                     &1
   \end{mat}
   \begin{mat}
     \frac{1+\sqrt{5}}{2}  &0   \\
     0                     &\frac{1-\sqrt{5}}{2}
   \end{mat}
   \begin{mat}
     \frac{1}{\sqrt{5}}     &-(\frac{1-\sqrt{5}}{2\sqrt{5}})  \\
     \frac{-1}{\sqrt{5}}    &\frac{1+\sqrt{5}}{2\sqrt{5}}       
   \end{mat}
\end{equation*} 
Dengan memperkenalkan vektor-vektor tersebut dan mengambil pangkat ke-$n$,
kita memperoleh
\begin{multline*}
  \colvec{F(n) \\ F(n-1)}
  =\begin{mat}[r]
    1  &1  \\
    1  &0
  \end{mat}^{n-1}
  \colvec{F(1) \\ F(0)}                                          \\
  =\begin{mat}
     \frac{1+\sqrt{5}}{2}  &\frac{1-\sqrt{5}}{2} \\
     1                     &1
   \end{mat}
   \begin{mat}
     \left(\frac{1+\sqrt{5}}{2}\right)^{n-1}  &0   \\
     0                                  &\left(\frac{1-\sqrt{5}}{2}\right)^{n-1}
   \end{mat}
   \begin{mat}
     \frac{1}{\sqrt{5}}     &-(\frac{1-\sqrt{5}}{2\sqrt{5}})  \\
     \frac{-1}{\sqrt{5}}    &\frac{1+\sqrt{5}}{2\sqrt{5}}       
   \end{mat}
  \colvec[r]{1 \\ 0}                                             
\end{multline*}
Perhitungannya kurang elok, tetapi tidak sulit.
\begin{align*}
  \colvec{F(n) \\ F(n-1)}
  &=\begin{mat}
     \frac{1+\sqrt{5}}{2}  &\frac{1-\sqrt{5}}{2} \\
     1                     &1
   \end{mat}
   \begin{mat}
     \left(\frac{1+\sqrt{5}}{2}\right)^{n-1}  &0   \\
     0                                  &\left(\frac{1-\sqrt{5}}{2}\right)^{n-1}
   \end{mat}
  \colvec{\frac{1}{\sqrt{5}} \\ -\frac{1}{\sqrt{5}}}       \\ 
  &=\frac{1}{\sqrt{5}}
    \begin{mat}
     \frac{1+\sqrt{5}}{2}  &\frac{1-\sqrt{5}}{2} \\
     1                     &1
   \end{mat}
  \colvec{\left(\frac{1+\sqrt{5}}{2}\right)^{n-1} \\ 
          -\left(\frac{1-\sqrt{5}}{2}\right)^{n-1}}    \\    
  &=\frac{1}{\sqrt{5}}
  \colvec{\left(\frac{1+\sqrt{5}}{2}\right)^{n}  
          -\left(\frac{1-\sqrt{5}}{2}\right)^{n}  \\
          \left(\frac{1+\sqrt{5}}{2}\right)^{n-1}  
          -\left(\frac{1-\sqrt{5}}{2}\right)^{n-1}}   
\end{align*}
Yang kita perlukan adalah komponen pertama.
\begin{equation*}
  F(n)=\frac{1}{\sqrt{5}}\left[\left(\frac{1+\sqrt{5}}{2}\right)^{n}
                               -\left(\frac{1-\sqrt{5}}{2}\right)^{n}\right]
\end{equation*}
Rumus ini memberikan nilai setiap suku barisan tanpa harus lebih dahulu
menentukan nilai-nilai antaranya.

Karena $(1-\sqrt{5})/2\approx -0.618$ bernilai mutlak kurang dari satu,
pangkat-pangkatnya menuju nol. Oleh karena itu, rumus $F(n)$ didominasi oleh
suku pertamanya.
Meskipun kita telah mengembangkan model dasar pertumbuhan populasi dengan
menambahkan masa tunda sebelum kesuburan dimulai, kita tetap memperoleh
fungsi yang secara asimtotik bersifat eksponensial.

Secara umum, suatu
\definend{relasi rekurensi linear homogen berkoefisien konstan dan berorde $k$}\index{rekurensi}\index{rekurensi linear!definisi}
berbentuk sebagai berikut.
\begin{equation*}
  f(n)=a_1f(n-1)+a_2f(n-2)+\dots+a_kf(n-k)
\end{equation*}
Relasi rekurensi ini homogen karena tidak mempunyai suku konstanta; dengan
kata lain, kita dapat menuliskannya kembali sebagai
$0=-f(n)+a_1f(n-1)+a_2f(n-2)+\dots+a_kf(n-k)$.
Relasi ini berorde~$k$ karena menggunakan $k$ suku sebelumnya untuk
menghitung $f(n)$.
Relasi tersebut, bersama dengan
kondisi awal\index{rekurensi!kondisi awal} yang memberikan nilai
$f(0)$, \ldots, $f(k-1)$, menentukan suatu barisan secara lengkap, sebab
kita dapat menghitung $f(n)$ dengan lebih dahulu menghitung $f(k)$,
$f(k+1)$, dan seterusnya.
Seperti pada kasus Fibonacci, kita akan mencari rumus yang menyelesaikan
rekurensi tersebut dan memberikan $f(n)$ secara langsung.

Misalkan $V$ adalah himpunan fungsi dengan domain
$\N=\set{0,1,2,\ldots}$ dan kodomain $\C$.
(Jika memudahkan, kadang-kadang kita menggunakan domain
$\Z^+=\set{1,2,\ldots}$.)
Himpunan ini merupakan ruang vektor dengan pengertian biasa untuk
penjumlahan dan perkalian skalar: $f+g$ adalah pemetaan
$x\mapsto f(x)+g(x)$ dan $cf$ adalah pemetaan $x\mapsto c\cdot f(x)$.

Jika kita mengesampingkan kondisi awal dan hanya memperhatikan
rekurensinya, mungkin terdapat banyak fungsi yang memenuhi relasi tersebut.
Sebagai contoh, rekurensi Fibonacci, yaitu bahwa setiap nilai setelah
nilai-nilai awal merupakan jumlah dua nilai sebelumnya, dipenuhi oleh fungsi
$L$ yang beberapa nilai pertamanya adalah $L(0)=2$, $L(1)=1$, $L(2)=3$,
$L(3)=4$, dan $L(4)=7$.

Tetapkan suatu relasi rekurensi linear homogen berorde~$k$, lalu perhatikan
subhimpunan $S$ yang terdiri atas fungsi-fungsi yang memenuhi relasi tersebut
(tanpa kondisi awal).
Himpunan $S$ ini adalah subruang dari $V$.
Himpunan ini tidak kosong karena, berdasarkan kehomogenan, fungsi nol adalah
sebuah solusi.
Himpunan ini tertutup terhadap penjumlahan karena, jika $f_1$ dan $f_2$
merupakan solusi, maka persamaan berikut berlaku.
\begin{multline*}
  -(f_1+f_2)(n)+a_1(f_1+f_2)(n-1)+\dots+a_k(f_1+f_2)(n-k) \\  
  \begin{aligned}
    &=(-f_1(n)+\dots+a_kf_1(n-k))          \\
    &\qquad\hbox{}+(-f_2(n)+\dots+a_kf_2(n-k))     \\
    &=0+0=0
  \end{aligned}
\end{multline*}
Himpunan ini juga tertutup terhadap perkalian skalar.
\begin{multline*}
  -(rf_1)(n)+a_1(rf_1)(n-1)+\dots+a_k(rf_1)(n-k) \\  
  \begin{aligned}
    &=r\cdot (-f_1(n)+\dots+a_kf_1(n-k))   \\
    &=r\cdot 0                                    \\
    &=0
  \end{aligned}
\end{multline*}
Kita dapat menentukan dimensi $S$.
Karena $k$ adalah orde rekurensi, perhatikan pemetaan berikut dari himpunan
fungsi~$S$ ke himpunan vektor dengan $k$ komponen.
\begin{equation*}
  f 
  \mapsto 
  \colvec{f(0) \\ f(1) \\ \vdots \\ f(k-1)}
\end{equation*}
\nearbyexercise{exer:SeqToRnLinMap} menunjukkan bahwa pemetaan ini linear.
Setiap solusi rekurensi ditentukan secara tunggal oleh $k$ kondisi awal,
sehingga pemetaan ini bersifat injektif dan surjektif.
Jadi, pemetaan ini merupakan isomorfisme dan $S$ berdimensi $k$.

Dengan demikian, kita dapat mendeskripsikan himpunan solusi relasi rekurensi
linear homogen berorde~$k$ dengan mencari suatu basis yang terdiri atas
$k$ fungsi bebas linear.
Untuk menghasilkan fungsi-fungsi itu, kita menyatakan rekurensi tersebut
dalam bentuk matriks.
\begin{equation*}
  \colvec{f(n) \\ f(n-1) \\ \vdots  \\ f(n-k+1)}
  =
  \begin{mat}
    a_1  &a_2  &a_3  &\ldots  &a_{k-1} &a_k  \\
    1    &0        &0        &\ldots  &0         &0        \\
    0    &1        &0                                      \\
    0    &0        &1                                      \\
    \vdots &\vdots &         &\ddots  &           &\vdots  \\
    0    &0        &0        &\ldots   &1          &0
  \end{mat}
  \colvec{f(n-1) \\ f(n-2) \\ \vdots  \\ f(n-k)}
\end{equation*}
Sebut matriks tersebut~$A$.
Kita ingin menentukan fungsi karakteristiknya, yaitu determinan
$A-\lambda I$.
Pola pada kasus $\nbyn{2}$
\begin{equation*}
  \begin{vmat}
    a_1-\lambda  &a_2 \\
    1               &-\lambda
  \end{vmat}
  =\lambda^2-a_1\lambda-a_2
\end{equation*}
dan pada kasus $\nbyn{3}$
\begin{equation*}
  \begin{vmat}
    a_1-\lambda  &a_2   &a_3  \\
    1               &-\lambda  &0        \\
    0               &1         &-\lambda
  \end{vmat}
  =-\lambda^3+a_1\lambda^2+a_2\lambda+a_3
\end{equation*}
mengarahkan kita untuk menduga, dan
\nearbyexercise{exer:CharEqnIsDeter} membuktikan, bahwa persamaan berikut
adalah persamaan karakteristiknya.
\begin{align*}
  0 &=
  \begin{vmat}
    a_1-\lambda &a_2  &a_3  &\ldots  &a_{k-1} &a_k  \\
    1              &-\lambda &0        &\ldots  &0         &0        \\
    0              &1        &-\lambda                                      \\
    0              &0        &1                                      \\
    \vdots         &\vdots &         &\ddots   &           &\vdots  \\
    0             &0        &0        &\ldots   &1          &-\lambda
  \end{vmat}                                                           \\
  &=\pm(-\lambda^k+a_1\lambda^{k-1}+a_2\lambda^{k-2}
       +\dots+a_{k-1}\lambda+a_k)
\end{align*}
Tanda $\pm$ tidak relevan untuk menentukan akar, sehingga kita abaikan.
Kita katakan bahwa polinomial
$-\lambda^k+a_1\lambda^{k-1}+a_2\lambda^{k-2}
       +\dots+a_{k-1}\lambda+a_k$
berasosiasi dengan relasi rekurensi tersebut.
\index{rekurensi!polinomial berasosiasi}\index{polinomial!berasosiasi dengan rekurensi}

Jika persamaan karakteristik tidak mempunyai akar berulang, matriks tersebut
dapat didiagonalkan dan, secara teori, kita dapat memperoleh rumus untuk
$f(n)$ seperti pada kasus Fibonacci.
Namun, karena kita mengetahui bahwa subruang solusi berdimensi~$k$, kita tidak
perlu melakukan perhitungan pendiagonalan asalkan kita dapat menunjukkan
$k$ fungsi berbeda yang bebas linear dan memenuhi relasi tersebut.

Jika $r_1$, $r_2$, \ldots, $r_k$ adalah akar-akar yang berbeda, perhatikan
fungsi-fungsi berupa pangkat akar tersebut, dari $f_{r_1}(n)=r_1^n$ hingga
$f_{r_k}(n)=r_k^n$.
\nearbyexercise{exer:SoltnsLinRecur} menunjukkan bahwa masing-masing fungsi
merupakan solusi rekurensi dan bahwa semuanya membentuk himpunan bebas linear.
% So given the homogeneous linear recurrence
% $f(n+1)=a_nf(n)+\dots+a_{n-k}f(n-k)$
% (that is, $0=-f(n+1)+a_nf(n)+\dots+a_{n-k}f(n-k)$)
% we consider the associated equation
% $0=-\lambda^k+a_n\lambda^{k-1}+\dots+a_{n-k+1}\lambda+a_{n-k}$.
Jadi, jika akar-akar polinomial yang berasosiasi itu berbeda, setiap solusi
relasi tersebut berbentuk
$f(n)=c_1r_1^n+c_2r_2^n+\dots+c_kr_k^n$ untuk beberapa skalar
$c_1, \dots, c_k$.
(Kasus akar berulang serupa, tetapi tidak akan kita bahas di sini; lihat buku
teks matematika diskret mana pun.)

Sekarang kita sertakan kondisi awal.
Gunakan kondisi-kondisi itu untuk menentukan $c_1$, \ldots, $c_k$.
Sebagai contoh, polinomial yang berasosiasi dengan relasi Fibonacci adalah
$-\lambda^2+\lambda+1$, yang akar-akarnya $r_1=(1+\sqrt{5})/2$ dan
$r_2=(1-\sqrt{5})/2$. Dengan demikian, setiap solusi rekurensi Fibonacci
berbentuk $f(n)=c_1((1+\sqrt{5})/2)^n+c_2((1-\sqrt{5})/2)^n$.
Gunakan kondisi awal Fibonacci untuk $n=0$ dan $n=1$,
\begin{equation*}
  \begin{linsys}{2}
    c_1               &+  &c_2               &=  &0  \\
    ((1+\sqrt{5})/2)c_1 &+  &((1-\sqrt{5})/2)c_2 &=  &1
  \end{linsys}
\end{equation*}
dan selesaikan untuk memperoleh $c_1=1/\sqrt{5}$ dan
$c_2=-1/\sqrt{5}$, seperti yang kita temukan di atas.

Sebagai penutup, kita tinjau kasus tak homogen, dengan relasi berbentuk
$f(n)=a_1f(n-1)+a_2f(n-2)+\dots+a_kf(n-k)+b$ untuk suatu $b$ taknol.
Kita hanya memerlukan sedikit penyesuaian untuk beralih dari kasus homogen.

Contoh klasik berikut memberikan ilustrasi.
Pada tahun 1883, Edouard Lucas mengajukan masalah
Menara Hanoi.\index{Menara Hanoi}
\begin{quotation}
  \noindent
  Di kuil besar di Benares, di bawah kubah yang menandai pusat dunia,
  terdapat sebuah lempeng kuningan tempat tertancap tiga jarum berlian,
  masing-masing setinggi satu hasta dan setebal tubuh seekor lebah.
  Pada saat penciptaan, Tuhan menaruh enam puluh empat cakram emas murni
  pada salah satu jarum ini; cakram terbesar terletak di atas lempeng
  kuningan, dan cakram-cakram lain makin mengecil hingga cakram teratas.
  Inilah Menara Brahma.
  Siang dan malam tanpa henti, para pendeta memindahkan cakram-cakram dari
  satu jarum berlian ke jarum lain menurut hukum Brahma yang tetap dan tak
  berubah. Hukum itu menetapkan bahwa pendeta yang bertugas tidak boleh
  memindahkan lebih dari satu cakram sekaligus dan harus menaruh cakram itu
  pada jarum sedemikian sehingga tidak ada cakram yang lebih kecil di
  bawahnya.
  Setelah keenam puluh empat cakram itu dipindahkan dari jarum tempat Tuhan
  menaruhnya pada saat penciptaan ke salah satu jarum lain, menara, kuil,
  dan para Brahmana akan sama-sama runtuh menjadi debu, dan dunia akan lenyap
  dengan satu gelegar guntur.
  (Terjemahan \cite{DeParville} dari \cite{Ball}.)
\end{quotation}
Kita mengesampingkan pertanyaan mengapa para pendeta tidak duduk beristirahat
agar dunia bertahan sedikit lebih lama. Sebagai gantinya, kita tanyakan berapa
banyak pemindahan cakram yang diperlukan.
Sebelum menangani masalah enam puluh empat cakram, kita akan meninjau masalah
untuk tiga cakram.

Pada awalnya, ketiga cakram berada pada jarum yang sama.
\begin{center}
  \includegraphics{jc/mp/ch5.6}
\end{center}
Setelah tiga langkah---memindahkan cakram kecil ke jarum terjauh, cakram
sedang ke jarum tengah, lalu cakram kecil ke jarum tengah---kita memperoleh
keadaan berikut.
\begin{center}
  \includegraphics{jc/mp/ch5.7}
\end{center}
Sekarang kita dapat memindahkan cakram besar ke jarum terjauh.
Untuk menyelesaikannya, kita ulangi proses tiga langkah pada dua cakram yang
lebih kecil, kali ini agar keduanya berakhir pada jarum ketiga, di atas
cakram besar.

Urutan langkah tersebut adalah yang terbaik yang dapat kita lakukan.
Untuk memindahkan cakram terbawah, setidaknya kita harus lebih dahulu
memindahkan cakram-cakram yang lebih kecil ke jarum tengah, lalu memindahkan
cakram besar, kemudian memindahkan semua cakram yang lebih kecil dari jarum
tengah ke jarum tujuan.
Karena jumlah minimum ini dapat dicapai, kita memperoleh rekurensi berikut.
\begin{equation*}
  T(n)=T(n-1)+1+T(n-1)=2T(n-1)+1 \qquad \text{dengan $T(1)=1$}
\end{equation*} 
Berikut adalah beberapa nilai pertama $T$.
\begin{center}
  \begin{tabular}{r|cccccccccc}
    \textit{cakram}~$n$             
       &1  &2  &3  &4  &5  &6     &7    &8    &9   &10  \\
    \hline
    \textit{pemindahan}~$T(n)$ 
       &1  &3  &7  &15  &31  &63  &127  &255  &511 &1023 
  \end{tabular}
\end{center}
Tentu saja, masing-masing bilangan ini kurang satu dari suatu pangkat dua.
Untuk menurunkan hal tersebut, tuliskan relasi semula sebagai
$-1=-T(n)+2T(n-1)$.
Perhatikan $0=-T(n)+2T(n-1)$, suatu rekurensi linear homogen berorde~$1$.
Polinomial yang berasosiasi dengannya adalah $-\lambda+2$, dengan satu-satunya
akar $r_1=2$.
Jadi, fungsi yang memenuhi relasi homogen berbentuk $c_12^n$.

Itulah solusi homogen.
Sekarang kita memerlukan suatu solusi partikular.
Karena relasi tak homogen $-1=-T(n)+2T(n-1)$ begitu sederhana, kita dapat
langsung melihat solusi partikular $T(n)=-1$.
Setiap solusi rekurensi $T(n)=2T(n-1)+1$ (tanpa kondisi awal) merupakan jumlah
solusi homogen dan solusi partikular, yaitu $c_12^n-1$.
Sekarang kondisi awal $T(1)=1$ memberikan $c_1=1$, dan kita telah memperoleh
rumus yang menghasilkan tabel tersebut:~masalah Menara Hanoi dengan $n$
cakram memerlukan $T(n)=2^n-1$ pemindahan.

Menentukan solusi partikular pada kasus yang lebih rumit, barangkali tidak
mengherankan, juga lebih rumit.
Sumber yang menyenangkan dan bermanfaat, tetapi menantang, adalah
\cite{GrahamKnuthPatashnik}.
Untuk pembahasan lebih lanjut tentang Menara Hanoi, lihat \cite{Ball},
\cite{Gardner57}, dan \cite{Hofstadter}.
Sejumlah kode komputer disajikan setelah latihan.

\begin{exercises}
  \item 
   Berapa bulan diperlukan hingga jumlah pasangan kelinci Fibonacci
   melampaui seribu?
   Sepuluh ribu?
   Satu juta?
   \begin{answer}
     Kita menggunakan rumus
     \begin{equation*}
       F(n)=\frac{1}{\sqrt{5}}\left[\left(\frac{1+\sqrt{5}}{2}\right)^{n}
                               -\left(\frac{1-\sqrt{5}}{2}\right)^{n}\right]
     \end{equation*}
     Seperti telah diamati sebelumnya, $(1+\sqrt{5})/2$ lebih besar dari satu,
     sedangkan $(1-\sqrt{5})/2$ bernilai mutlak kurang dari satu.
\begin{lstlisting}
sage: phi = (1+5^(0.5))/2
sage: psi = (1-5^(0.5))/2
sage: phi
1.61803398874989
sage: psi
-0.618033988749895       
\end{lstlisting}
     Jadi, nilai ungkapan tersebut didominasi oleh suku pertama.
     Menyelesaikan $1000=(1/\sqrt{5})\cdot ((1+\sqrt{5})/2)^n$
     memberikan hasil berikut.
\begin{lstlisting}
sage: a = ln(1000*5^(0.5))/ln(phi)
sage: a
16.0271918385296
sage: psi^(17)
-0.000280033582072583       
\end{lstlisting}
   Jadi, pada pangkat ketujuh belas, sumbangan suku kedua tidak cukup besar
   untuk mengubah hasil pembulatan.
   Untuk perhitungan sepuluh ribu dan satu juta, situasinya bahkan lebih
   ekstrem.
\begin{lstlisting}
sage: b = ln(10000*5^(0.5))/ln(phi)
sage: b
20.8121638053112
sage: c = ln(1000000*5^(0.5))/ln(phi)
sage: c
30.3821077388746    
\end{lstlisting}
   Jawaban untuk kedua kasus ini adalah $21$ dan $31$.
   \end{answer}
  \item 
   Selesaikan setiap relasi rekurensi linear homogen berikut.
   \begin{exparts}
     \partsitem $f(n)=5f(n-1)-6f(n-2)$   
     \partsitem $f(n)=4f(n-2)$   
     \partsitem $f(n)=5f(n-1)-2f(n-2)-8f(n-3)$   
   \end{exparts}
   \begin{answer}
    \begin{exparts}
      \partsitem 
        Kita menyatakan relasi tersebut dalam bentuk matriks.
        \begin{equation*}
          \colvec{f(n) \\ f(n-1)}
          =
          \begin{mat}[r]
            5  &-6  \\
            1  &0
          \end{mat}
          \colvec{f(n-1) \\ f(n-2)}
        \end{equation*}
        Persamaan karakteristik matriks tersebut,
        \begin{equation*}
          \begin{vmat}
            5-\lambda &-6       \\
            1         &-\lambda
          \end{vmat}
          =\lambda^2-5\lambda+6 
        \end{equation*}
        mempunyai akar $2$ dan $3$.
        Setiap fungsi berbentuk
        $f(n)=c_12^n+c_23^n$
        memenuhi rekurensi tersebut.
      \partsitem 
        Bentuk matriks relasi tersebut adalah
        \begin{equation*}
          \colvec{f(n) \\ f(n-1)}
          =
          \begin{mat}[r]
            0  &4  \\
            1  &0  
          \end{mat}
          \colvec{f(n-1) \\ f(n-2)}
        \end{equation*}
        dan persamaan karakteristiknya,
        \begin{equation*}
          \begin{vmat}
            -\lambda &4  \\
            1        &-\lambda
          \end{vmat}
          =\lambda^2-4=(\lambda-2)(\lambda+2)
        \end{equation*}
        mempunyai dua akar, yaitu $2$ dan $-2$.
        Setiap fungsi berbentuk
        $f(n)=c_12^n+c_2(-2)^n$
        memenuhi rekurensi ini.
      \partsitem 
        Dalam bentuk matriks, relasi
        \begin{equation*}
          \colvec{f(n) \\ f(n-1) \\ f(n-2)}
          =
          \begin{mat}[r]
            5  &-2  &-8  \\
            1  &0  &0  \\
            0  &1  &0
          \end{mat}
          \colvec{f(n-1) \\ f(n-2) \\ f(n-3)}
        \end{equation*}
        mempunyai persamaan karakteristik dengan akar $-1$, $2$, dan $4$.
        Setiap kombinasi berbentuk
        $c_1(-1)^n+c_22^n+c_34^n$ menyelesaikan rekurensi tersebut.
    \end{exparts} 
   \end{answer}
  \item \label{exer:SolvePartRecurSoltn} 
    Berikan rumus untuk relasi-relasi pada latihan sebelumnya dengan kondisi
    awal berikut.
    \begin{exparts}
      \partsitem $f(0)=1$, $f(1)=1$
      \partsitem $f(0)=0$, $f(1)=1$
      \partsitem $f(0)=1$, $f(1)=1$, $f(2)=3$.      
    \end{exparts}
    \begin{answer}
      \begin{exparts}
        \partsitem Solusi rekurensi homogen tersebut adalah
          $f(n)=c_12^n+c_23^n$.
          Menyubstitusikan $f(0)=1$ dan $f(1)=1$ memberikan sistem linear
          berikut.
          \begin{equation*}
            \begin{linsys}{2}
              c_1 &+ &c_2   &= &1 \\
             2c_1 &+ &3c_2  &= &1
            \end{linsys}
          \end{equation*}
          Dengan pengamatan langsung, kita memperoleh $c_1=2$ dan $c_2=-1$.
        \partsitem
          Solusi rekurensi homogen tersebut adalah
          $c_12^n+c_2(-2)^n$.
          Kondisi awal memberikan sistem linear berikut.
          \begin{equation*}
            \begin{linsys}{2}
              c_1 &+ &c_2  &= &0  \\
             2c_1 &- &2c_2 &= &1
            \end{linsys}
          \end{equation*}
          Solusinya adalah $c_1=1/4$, $c_2=-1/4$.
        \partsitem
          Rekurensi homogen tersebut mempunyai solusi
          $f(n)=c_1(-1)^n+c_22^n+c_34^n$.
          Dengan kondisi awal, kita memperoleh sistem linear berikut.
          \begin{equation*}
            \begin{linsys}{3}
              c_1 &+  &c_2   &+  &c_3  &= &1 \\
             -c_1 &+  &2c_2  &+  &4c_3  &= &1 \\
              c_1 &+  &4c_2  &+  &16c_3  &= &3
            \end{linsys}
          \end{equation*}
          Solusinya adalah $c_1=1/3$, $c_2=2/3$, $c_3=0$.
      \end{exparts}
    \end{answer}
  \item \label{exer:SeqToRnLinMap}
    Periksa bahwa isomorfisme yang diberikan antara $S$ dan $\C^k$ merupakan
    pemetaan linear.
    % We argue above that this map is one-to-one.
    % What is its inverse?
    \begin{answer}
      Tetapkan suatu rekurensi linear homogen berorde~$k$.
      Misalkan $S$ adalah himpunan fungsi $\map{f}{\N}{\C}$ yang memenuhi
      rekurensi tersebut.
      Perhatikan fungsi $\map{\Phi}{S}{\C^k}$ berikut.
      \begin{equation*}
        f\mapsunder{\Phi}\colvec{f(0) \\ f(1) \\ \vdots \\ f(k-1)}
      \end{equation*}
      Perhitungan berikut menunjukkan linearitasnya.
      \begin{multline*}
        \Phi(a\cdot f_1+b\cdot f_2)
        =
        \colvec{af_1(0)+bf_2(0) \\ \vdots \\ af_1(k-1)+bf_2(k-1)}     \\
        =
        a\colvec{f_1(0) \\ \vdots \\ f_1(k-1)}
        +
        b\colvec{f_2(0) \\ \vdots \\ f_2(k-1)}
        =a\cdot \Phi(f_1)+b\cdot \Phi(f_2)
      \end{multline*}
    \end{answer}
  \item \label{exer:CharEqnIsDeter}
    Tunjukkan bahwa persamaan karakteristik matriks tersebut adalah seperti
    yang dinyatakan, yaitu polinomial yang berasosiasi dengan relasi itu.
    (\textit{Petunjuk:} ekspansi sepanjang kolom terakhir dan penggunaan
    induksi akan berhasil.)
    \begin{answer}
      Kita menggunakan petunjuk tersebut untuk membuktikannya.
      \begin{align*}
        0 &=
        \begin{vmat}
          a_1-\lambda &a_2  &a_3  &\ldots  &a_{k-1} &a_k  \\
          1              &-\lambda &0        &\ldots  &0         &0        \\
          0              &1        &-\lambda                               \\
          0              &0        &1                                      \\
          \vdots         &\vdots &         &\ddots   &           &\vdots  \\
          0             &0        &0        &\ldots   &1          &-\lambda
        \end{vmat}                                                           \\
        &=(-1)^{k-1}(-\lambda^k+a_1\lambda^{k-1}+a_2\lambda^{k-2}       
               +\dots+a_{k-1}\lambda+a_k)                      \\
        &=\pm(-\lambda^k+a_1\lambda^{k-1}+a_2\lambda^{k-2}       
               +\dots+a_{k-1}\lambda+a_k)
      \end{align*}
      Langkah basis bersifat langsung.
      Untuk langkah induksi, ekspansi sepanjang kolom terakhir memberikan
      dua suku taknol.
      \begin{equation*}
        (-1)^{k-1}a_k\cdot 1
        -\lambda\cdot
        \begin{vmat}
          a_1-\lambda &a_2  &a_3  &\ldots  &a_{k-1}  \\
          1              &-\lambda &0        &\ldots  &0       \\
          0              &1        &-\lambda                               \\
          0              &0        &1                                      \\
          \vdots         &\vdots &         &\ddots   &        \\
          0             &0        &0        &\ldots   &1     
        \end{vmat}                                                          
      \end{equation*}
      (Matriks tersebut persegi, sehingga tanda di depan $-\lambda$ adalah
      $(-1)^{\text{genap}}$.)
      Penerapan hipotesis induksi memberikan hasil yang diinginkan.
      \begin{multline*}
        =(-1)^{k-1}a_k\cdot 1                     \\
        -\lambda\cdot
         (-1)^{k-2}(-\lambda^{k-1}+a_1\lambda^{k-2}+a_2\lambda^{k-3}       
               +\dots+a_{k-1}\lambda^0)
      \end{multline*}
    \end{answer}
  \item \label{exer:SoltnsLinRecur}
    Diberikan suatu relasi rekurensi linear homogen
    $f(n)=a_1f(n-1)+\dots+a_kf(n-k)$, misalkan $r_1$, \ldots, $r_k$
    adalah akar-akar polinomial yang berasosiasi dengannya.
    Buktikan bahwa setiap fungsi
         $f_{r_i}(n)=r_i^n$
         memenuhi rekurensi tersebut (tanpa kondisi awal).
    \begin{answer}
       Hasil ini diperoleh dengan induksi langsung pada $n$.
    \end{answer}
  \item 
    (Soal ini merujuk pada nilai $T(64)=18,446,744,073,709,551,615$ yang
    diberikan dalam kode komputer di bawah.)
    Jika satu cakram dipindahkan setiap detik, berapa tahun yang diperlukan
    para pendeta di Menara Hanoi untuk menyelesaikan tugas tersebut?
    \begin{answer}
      \Sage{} menyatakan bahwa kita masih aman.
\begin{lstlisting}
sage: T64 = 18446744073709551615
sage: T64_days = T64/(60*60*24)
sage: T64_days
1229782938247303441/5760
sage: T64_years = T64_days/365.25
sage: T64_years
5.84542046090626e11        
sage: age_of_universe = 13.8e9
sage: T64_years/age_of_universe
42.3581192819294
\end{lstlisting}
    \end{answer}

\end{exercises}

\announcecomputercode
Kode ini menghasilkan beberapa nilai pertama suatu fungsi yang didefinisikan
oleh rekurensi dan kondisi awal.
Kode tersebut ditulis dalam dialek Scheme dari LISP, khususnya
\cite{ChickenScheme}.

Setelah memuat ekstensi yang mencegah komputer beralih ke bilangan titik
mengambang ketika bilangan bulat menjadi besar, fungsi Menara Hanoi dapat
ditulis dengan mudah.
\begin{lstlisting}
(require-extension numbers)

(define (tower-of-hanoi-moves n) 
    (if (= n 1)
       1
       (+ (* (tower-of-hanoi-moves (- n 1)) 
              2) 
           1) ) )

; Dua fungsi bantu
(define (first-few-outputs proc n)
    (first-few-outputs-aux proc n '()) )

(define (first-few-outputs-aux proc n lst)
    (if (< n 1)
    lst 
    (first-few-outputs-aux proc (- n 1) (cons (proc n) lst)) ) )
\end{lstlisting}
\noindent (Bagi pembaca yang belum terbiasa dengan kode rekursif:~untuk
menghitung $T(64)$, komputer hendak menghitung $2*T(63)+1$, yang mengharuskan
komputer menghitung $T(63)$.
Komputer menyisihkan operasi `kali $2$' dan `tambah $1$' sejenak.
Komputer menghitung $T(63)$ menggunakan potongan kode yang sama (itulah arti
`rekursif'), dan untuk melakukannya komputer hendak menghitung
$2*T(62)+1$.
Proses ini berlanjut hingga, setelah $63$ langkah, komputer mencoba menghitung
$T(1)$.
Komputer lalu mengembalikan $T(1)=1$, yang memungkinkan perhitungan $T(2)$
berlanjut, dan seterusnya, hingga perhitungan awal $T(64)$ selesai.)

Fungsi-fungsi pembantu memberikan tabel beberapa nilai pertama.
% (Some language notes:~\texttt{'()} is the empty list
% and \texttt{cons} pushes something onto the start of a 
% list.
% Note that, in the last line, the function \texttt{proc}
% is called on argument \texttt{n}.)
Berikut adalah sesi pada baris perintah.
\begin{lstlisting}
#;1> (load "hanoi.scm")
; loading hanoi.scm ...
; loading /var/lib//chicken/6/numbers.import.so ...
; loading /var/lib//chicken/6/chicken.import.so ...
; loading /var/lib//chicken/6/foreign.import.so ...
; loading /var/lib//chicken/6/numbers.so ...
#;2> (tower-of-hanoi-moves 64)
18446744073709551615
#;3> (first-few-outputs tower-of-hanoi-moves 64)
(1 3 7 15 31 63 127 255 511 1023 2047 4095 8191 16383 32767 65535 131071 262143 524287 1048575 
2097151 4194303 8388607 16777215 33554431 67108863 134217727 268435455 536870911 1073741823 
2147483647 4294967295 8589934591 17179869183 34359738367 68719476735 137438953471 274877906943 
549755813887 1099511627775 2199023255551 4398046511103 8796093022207 17592186044415 
35184372088831 70368744177663 140737488355327 281474976710655 562949953421311 1125899906842623 
2251799813685247 4503599627370495 9007199254740991 18014398509481983 36028797018963967 
72057594037927935 144115188075855871 288230376151711743 576460752303423487 1152921504606846975 
2305843009213693951 4611686018427387903 9223372036854775807 18446744073709551615)
\end{lstlisting}
\noindent Ini adalah daftar dari $T(1)$ hingga $T(64)$ (sesi tersebut telah
disunting dengan menambahkan pemisah baris agar mudah dibaca).
\index{relasi rekurensi|)}
\index{rekurensi linear|)}
\endinput
